Stack modifications are based on the
abstraction of stack locations, not the contents of these locations.
All stack offsets are with respect to the original
stack frame, even if \code{BPatch\_fuction::addMods} is called
multiple times for a single function.
Only implemented on x86 and x86-64.

\apidesc{
  This class defines modifications to the stack frame layout of a function.
}

\subsubsection{Insert}

\begin{apient}
  Insert(int low, int high)
\end{apient}

\apidesc{
  This constructor creates a stack modification that inserts stack
  space in the range
  [low, high), where low and high are stack offsets.
  \code{BPatch\_function::addMods}  will find this modification
  unsafe if any instructions in the function access memory that will
  be non-contiguous after [low,high) is inserted.
}

\subsubsection{Remove}

\begin{apient}
  Remove(int low, int high)
\end{apient}

\apidesc{
  This constructor creates a stack modification that removes stack
  space in the range
  [low, high), where low and high are stack offsets.
  \code{BPatch\_function::addMods}  will find this modification
  unsafe if any instructions in the function access stack memory in
  [low,high).
}

\subsubsection{Move}

\begin{apient}
  Move(int sLow, int sHigh, int dLow)
\end{apient}

\apidesc{
  This constructor creates a stack modification that moves stack
  space \code{[sLow, sHigh)} to \code{[dLow, dLow+(sHigh-sLow))}.
  \code{BPatch\_function::addMods}  will find this modification unsafe if
  \code{Insert(dLow, dLow+(sHigh-sLow))} or \code{Remove(sLow,
  sHigh)} are unsafe.
}

\subsubsection{Canary}

\begin{apient}
  Canary()
  Canary(BPatch_function* failFunc)
\end{apient}

\apidesc{
  implemented on Linux, GCC only
  
  This constructor creates a stack modification that inserts a stack
  canary at function entry and a corresponding canary
  check at function exit(s).
  This uses the same canary as GCC's \code{-fstack-protector}. If the
  canary check at function exit fails,
  failFunc is called. failFunc must be non-returning and take no
  arguments. If no failFunc is provided,
  \code{\_\_stack\_chk\_fail} from libc is called; libc must be open
  in the corresponding \code{BPatch\_addressSpace.}
  This modification will have no effect on functions in which the
  entry and exit point(s) are the same.
  \code{BPatch\_function::addMods}  will find this modification
  unsafe if another Canary has already been added to the function.
  Note, however, that this modification can be applied to code
  compiled with \code{-fstack-protector}.
}

\subsubsection{Randomize}

\begin{apient}
  Randomize()
  Randomize(int seed)
\end{apient}

\apidesc{
  This constructor creates a stack modification that rearranges the
  stack-stored local variables of a function. This
  modification requires symbol information (e.g., DWARF), and only
  local variables specified by the symbols will be
  randomized. If DyninstAPI finds a stack access that is not
  consistent with a symbol-specified local, that local will
  not be randomized. Contiguous ranges of local variables are
  randomized; if there are two or more contiguous ranges of
  locals within the stack frame, each is randomized separately. More
  than one local variable is required for
  randomization.
  \code{BPatch\_function::addMods}  will return false if Randomize is
  added to a function without local variable information,
  without local variables on the stack, or with only a single local variable.
  srand is used to generate a new ordering of local variables; if
  seed is provided, this value is provided to srand as its
  seed.
  \code{BPatch\_function::addMods}  will find this modification
  unsafe if any other modifications have been applied.
}
